% Options for packages loaded elsewhere
% Options for packages loaded elsewhere
\PassOptionsToPackage{unicode}{hyperref}
\PassOptionsToPackage{hyphens}{url}
\PassOptionsToPackage{dvipsnames,svgnames,x11names}{xcolor}
%
\documentclass[
  12pt,
  letterpaper,
  DIV=11,
  numbers=noendperiod]{scrartcl}
\usepackage{xcolor}
\usepackage{amsmath,amssymb}
\setcounter{secnumdepth}{5}
\usepackage{iftex}
\ifPDFTeX
  \usepackage[T1]{fontenc}
  \usepackage[utf8]{inputenc}
  \usepackage{textcomp} % provide euro and other symbols
\else % if luatex or xetex
  \usepackage{unicode-math} % this also loads fontspec
  \defaultfontfeatures{Scale=MatchLowercase}
  \defaultfontfeatures[\rmfamily]{Ligatures=TeX,Scale=1}
\fi
\usepackage{lmodern}
\ifPDFTeX\else
  % xetex/luatex font selection
  \setmainfont[]{Times New Roman}
\fi
% Use upquote if available, for straight quotes in verbatim environments
\IfFileExists{upquote.sty}{\usepackage{upquote}}{}
\IfFileExists{microtype.sty}{% use microtype if available
  \usepackage[]{microtype}
  \UseMicrotypeSet[protrusion]{basicmath} % disable protrusion for tt fonts
}{}
\makeatletter
\@ifundefined{KOMAClassName}{% if non-KOMA class
  \IfFileExists{parskip.sty}{%
    \usepackage{parskip}
  }{% else
    \setlength{\parindent}{0pt}
    \setlength{\parskip}{6pt plus 2pt minus 1pt}}
}{% if KOMA class
  \KOMAoptions{parskip=half}}
\makeatother
% Make \paragraph and \subparagraph free-standing
\makeatletter
\ifx\paragraph\undefined\else
  \let\oldparagraph\paragraph
  \renewcommand{\paragraph}{
    \@ifstar
      \xxxParagraphStar
      \xxxParagraphNoStar
  }
  \newcommand{\xxxParagraphStar}[1]{\oldparagraph*{#1}\mbox{}}
  \newcommand{\xxxParagraphNoStar}[1]{\oldparagraph{#1}\mbox{}}
\fi
\ifx\subparagraph\undefined\else
  \let\oldsubparagraph\subparagraph
  \renewcommand{\subparagraph}{
    \@ifstar
      \xxxSubParagraphStar
      \xxxSubParagraphNoStar
  }
  \newcommand{\xxxSubParagraphStar}[1]{\oldsubparagraph*{#1}\mbox{}}
  \newcommand{\xxxSubParagraphNoStar}[1]{\oldsubparagraph{#1}\mbox{}}
\fi
\makeatother

\usepackage{color}
\usepackage{fancyvrb}
\newcommand{\VerbBar}{|}
\newcommand{\VERB}{\Verb[commandchars=\\\{\}]}
\DefineVerbatimEnvironment{Highlighting}{Verbatim}{commandchars=\\\{\}}
% Add ',fontsize=\small' for more characters per line
\usepackage{framed}
\definecolor{shadecolor}{RGB}{241,243,245}
\newenvironment{Shaded}{\begin{snugshade}}{\end{snugshade}}
\newcommand{\AlertTok}[1]{\textcolor[rgb]{0.68,0.00,0.00}{#1}}
\newcommand{\AnnotationTok}[1]{\textcolor[rgb]{0.37,0.37,0.37}{#1}}
\newcommand{\AttributeTok}[1]{\textcolor[rgb]{0.40,0.45,0.13}{#1}}
\newcommand{\BaseNTok}[1]{\textcolor[rgb]{0.68,0.00,0.00}{#1}}
\newcommand{\BuiltInTok}[1]{\textcolor[rgb]{0.00,0.23,0.31}{#1}}
\newcommand{\CharTok}[1]{\textcolor[rgb]{0.13,0.47,0.30}{#1}}
\newcommand{\CommentTok}[1]{\textcolor[rgb]{0.37,0.37,0.37}{#1}}
\newcommand{\CommentVarTok}[1]{\textcolor[rgb]{0.37,0.37,0.37}{\textit{#1}}}
\newcommand{\ConstantTok}[1]{\textcolor[rgb]{0.56,0.35,0.01}{#1}}
\newcommand{\ControlFlowTok}[1]{\textcolor[rgb]{0.00,0.23,0.31}{\textbf{#1}}}
\newcommand{\DataTypeTok}[1]{\textcolor[rgb]{0.68,0.00,0.00}{#1}}
\newcommand{\DecValTok}[1]{\textcolor[rgb]{0.68,0.00,0.00}{#1}}
\newcommand{\DocumentationTok}[1]{\textcolor[rgb]{0.37,0.37,0.37}{\textit{#1}}}
\newcommand{\ErrorTok}[1]{\textcolor[rgb]{0.68,0.00,0.00}{#1}}
\newcommand{\ExtensionTok}[1]{\textcolor[rgb]{0.00,0.23,0.31}{#1}}
\newcommand{\FloatTok}[1]{\textcolor[rgb]{0.68,0.00,0.00}{#1}}
\newcommand{\FunctionTok}[1]{\textcolor[rgb]{0.28,0.35,0.67}{#1}}
\newcommand{\ImportTok}[1]{\textcolor[rgb]{0.00,0.46,0.62}{#1}}
\newcommand{\InformationTok}[1]{\textcolor[rgb]{0.37,0.37,0.37}{#1}}
\newcommand{\KeywordTok}[1]{\textcolor[rgb]{0.00,0.23,0.31}{\textbf{#1}}}
\newcommand{\NormalTok}[1]{\textcolor[rgb]{0.00,0.23,0.31}{#1}}
\newcommand{\OperatorTok}[1]{\textcolor[rgb]{0.37,0.37,0.37}{#1}}
\newcommand{\OtherTok}[1]{\textcolor[rgb]{0.00,0.23,0.31}{#1}}
\newcommand{\PreprocessorTok}[1]{\textcolor[rgb]{0.68,0.00,0.00}{#1}}
\newcommand{\RegionMarkerTok}[1]{\textcolor[rgb]{0.00,0.23,0.31}{#1}}
\newcommand{\SpecialCharTok}[1]{\textcolor[rgb]{0.37,0.37,0.37}{#1}}
\newcommand{\SpecialStringTok}[1]{\textcolor[rgb]{0.13,0.47,0.30}{#1}}
\newcommand{\StringTok}[1]{\textcolor[rgb]{0.13,0.47,0.30}{#1}}
\newcommand{\VariableTok}[1]{\textcolor[rgb]{0.07,0.07,0.07}{#1}}
\newcommand{\VerbatimStringTok}[1]{\textcolor[rgb]{0.13,0.47,0.30}{#1}}
\newcommand{\WarningTok}[1]{\textcolor[rgb]{0.37,0.37,0.37}{\textit{#1}}}

\usepackage{longtable,booktabs,array}
\newcounter{none} % for unnumbered tables
\usepackage{calc} % for calculating minipage widths
% Correct order of tables after \paragraph or \subparagraph
\usepackage{etoolbox}
\makeatletter
\patchcmd\longtable{\par}{\if@noskipsec\mbox{}\fi\par}{}{}
\makeatother
% Allow footnotes in longtable head/foot
\IfFileExists{footnotehyper.sty}{\usepackage{footnotehyper}}{\usepackage{footnote}}
\makesavenoteenv{longtable}
\usepackage{graphicx}
\makeatletter
\newsavebox\pandoc@box
\newcommand*\pandocbounded[1]{% scales image to fit in text height/width
  \sbox\pandoc@box{#1}%
  \Gscale@div\@tempa{\textheight}{\dimexpr\ht\pandoc@box+\dp\pandoc@box\relax}%
  \Gscale@div\@tempb{\linewidth}{\wd\pandoc@box}%
  \ifdim\@tempb\p@<\@tempa\p@\let\@tempa\@tempb\fi% select the smaller of both
  \ifdim\@tempa\p@<\p@\scalebox{\@tempa}{\usebox\pandoc@box}%
  \else\usebox{\pandoc@box}%
  \fi%
}
% Set default figure placement to htbp
\def\fps@figure{htbp}
\makeatother


% definitions for citeproc citations
\NewDocumentCommand\citeproctext{}{}
\NewDocumentCommand\citeproc{mm}{%
  \begingroup\def\citeproctext{#2}\cite{#1}\endgroup}
\makeatletter
 % allow citations to break across lines
 \let\@cite@ofmt\@firstofone
 % avoid brackets around text for \cite:
 \def\@biblabel#1{}
 \def\@cite#1#2{{#1\if@tempswa , #2\fi}}
\makeatother
\newlength{\cslhangindent}
\setlength{\cslhangindent}{1.5em}
\newlength{\csllabelwidth}
\setlength{\csllabelwidth}{3em}
\newenvironment{CSLReferences}[2] % #1 hanging-indent, #2 entry-spacing
 {\begin{list}{}{%
  \setlength{\itemindent}{0pt}
  \setlength{\leftmargin}{0pt}
  \setlength{\parsep}{0pt}
  % turn on hanging indent if param 1 is 1
  \ifodd #1
   \setlength{\leftmargin}{\cslhangindent}
   \setlength{\itemindent}{-1\cslhangindent}
  \fi
  % set entry spacing
  \setlength{\itemsep}{#2\baselineskip}}}
 {\end{list}}
\usepackage{calc}
\newcommand{\CSLBlock}[1]{\hfill\break\parbox[t]{\linewidth}{\strut\ignorespaces#1\strut}}
\newcommand{\CSLLeftMargin}[1]{\parbox[t]{\csllabelwidth}{\strut#1\strut}}
\newcommand{\CSLRightInline}[1]{\parbox[t]{\linewidth - \csllabelwidth}{\strut#1\strut}}
\newcommand{\CSLIndent}[1]{\hspace{\cslhangindent}#1}



\setlength{\emergencystretch}{3em} % prevent overfull lines

\providecommand{\tightlist}{%
  \setlength{\itemsep}{0pt}\setlength{\parskip}{0pt}}



 


\usepackage{tcolorbox}
\usepackage{amssymb}
\usepackage{yfonts}
\usepackage{bm}


\newtcolorbox{greybox}{
  colback=white,
  colframe=blue,
  coltext=black,
  boxsep=5pt,
  arc=4pt}
  
\newcommand{\sectionbreak}{\clearpage}

 
\newcommand{\ds}[4]{\sum_{{#1}=1}^{#3}\sum_{{#2}=1}^{#4}}
\newcommand{\us}[3]{\mathop{\sum\sum}_{1\leq{#2}<{#1}\leq{#3}}}

\newcommand{\ol}[1]{\overline{#1}}
\newcommand{\ul}[1]{\underline{#1}}

\newcommand{\amin}[1]{\mathop{\text{argmin}}_{#1}}
\newcommand{\amax}[1]{\mathop{\text{argmax}}_{#1}}

\newcommand{\ci}{\perp\!\!\!\perp}

\newcommand{\mc}[1]{\mathcal{#1}}
\newcommand{\mb}[1]{\mathbb{#1}}
\newcommand{\mf}[1]{\mathfrak{#1}}

\newcommand{\eps}{\epsilon}
\newcommand{\lbd}{\lambda}
\newcommand{\alp}{\alpha}
\newcommand{\df}{=:}
\newcommand{\am}[1]{\mathop{\text{argmin}}_{#1}}
\newcommand{\ls}[2]{\mathop{\sum\sum}_{#1}^{#2}}
\newcommand{\ijs}{\mathop{\sum\sum}_{1\leq i<j\leq n}}
\newcommand{\jis}{\mathop{\sum\sum}_{1\leq j<i\leq n}}
\newcommand{\sij}{\sum_{i=1}^n\sum_{j=1}^n}
	
\KOMAoption{captions}{tableheading}
\makeatletter
\@ifpackageloaded{caption}{}{\usepackage{caption}}
\AtBeginDocument{%
\ifdefined\contentsname
  \renewcommand*\contentsname{Table of contents}
\else
  \newcommand\contentsname{Table of contents}
\fi
\ifdefined\listfigurename
  \renewcommand*\listfigurename{List of Figures}
\else
  \newcommand\listfigurename{List of Figures}
\fi
\ifdefined\listtablename
  \renewcommand*\listtablename{List of Tables}
\else
  \newcommand\listtablename{List of Tables}
\fi
\ifdefined\figurename
  \renewcommand*\figurename{Figure}
\else
  \newcommand\figurename{Figure}
\fi
\ifdefined\tablename
  \renewcommand*\tablename{Table}
\else
  \newcommand\tablename{Table}
\fi
}
\@ifpackageloaded{float}{}{\usepackage{float}}
\floatstyle{ruled}
\@ifundefined{c@chapter}{\newfloat{codelisting}{h}{lop}}{\newfloat{codelisting}{h}{lop}[chapter]}
\floatname{codelisting}{Listing}
\newcommand*\listoflistings{\listof{codelisting}{List of Listings}}
\makeatother
\makeatletter
\makeatother
\makeatletter
\@ifpackageloaded{caption}{}{\usepackage{caption}}
\@ifpackageloaded{subcaption}{}{\usepackage{subcaption}}
\makeatother
\usepackage{bookmark}
\IfFileExists{xurl.sty}{\usepackage{xurl}}{} % add URL line breaks if available
\urlstyle{same}
\hypersetup{
  pdftitle={Convergence of SMACOF},
  pdfauthor={Jan de Leeuw},
  colorlinks=true,
  linkcolor={blue},
  filecolor={Maroon},
  citecolor={Blue},
  urlcolor={Blue},
  pdfcreator={LaTeX via pandoc}}


\title{Convergence of SMACOF}
\author{Jan de Leeuw}
\date{July 9, 2026}
\begin{document}
\maketitle
\begin{abstract}
TBD
\end{abstract}

\renewcommand*\contentsname{Table of contents}
{
\setcounter{tocdepth}{3}
\tableofcontents
}

\section{Introduction}\label{introduction}

The components of a (Euclidean, metric, least squares, r-dimensional)
multidimensional scaling (MDS) problem are

\begin{itemize}
\tightlist
\item
  a symmetric, non-negative, and
  hollow\footnote{A hollow matrix has a zero diagonal.} matrix
  \(\Delta\) of order \(n\) with \emph{dissimilarities},
\item
  a symmetric, non-negative, hollow, and
  irreducible\footnote{A matrix is irreducible if there is no permutation of rows and columns that makes it
  block-diagonal.} matrix \(W\) of order \(n\) with \emph{weights},
\item
  a \emph{basis} of \(m\) linearly independent \emph{configurations}
  \(Y_1,\cdots,Y_m\) of dimension \(n\times r\).
\end{itemize}

We define \emph{configuration space} as the set of all linear
combinations of the \(Y_s\), i.e. the set of all \(n\times r\) matrices
\(X\) for which there is an m-element vector \(\beta\) such that
\(\smash{X(\beta)=\sum_{s=1}^m\beta_sY_s}\).
Define\footnote{The symbol $:=$ is used for definitions.}

\[
d_{ij}(\beta):=\sqrt{(x_{i}(\beta)-x_{j}(\beta))'(x_{i}(\beta)-x_{j}(\beta))},
\] the Euclidean distance between rows \(i\) and \(j\) of \(X(\beta)\).
We minimize the real valued loss function \begin{equation}
\sigma(\beta):=\frac12\mathop{\sum}_{1\leq i<j\leq n} w_{ij}(\delta_{ij}-d_{ij}(\beta))^2,
\label{eq-stressx}
\end{equation} over \(\beta\) in the \emph{coefficient space}
\(\mathbb{R}^m\). We call \(\sigma(\beta)\) the \emph{stress} of
\emph{configuration} \(X(\beta)\) (Kruskal
(\citeproc{ref-kruskal_64a}{1964a}), Kruskal
(\citeproc{ref-kruskal_64b}{1964b})).

e parametrization \eqref{eq-alpha} is merely a linear change of
coordinatesIf \(m=np\) then the alternativ. Stress from
\eqref{eq-stressx} is the special case in which each of the \(np\)
matrices \(Y_s\) have a single element equal to one with all other
elements equal to zero. If \(m<np\) the constraint can be used to impose
restrictions on the configuration (Bentler and Weeks
(\citeproc{ref-bentler_weeks_78}{1978}), De Leeuw and Heiser
(\citeproc{ref-deleeuw_heiser_C_80}{1980})).

\sectionbreak

\section{Notation}\label{notation}

\(s,t=1,\cdots,m\)

\(p,q=1,\cdots,r\)

\(i,j,k,l=1,\cdots,n\)

First some convenient notation in configuration space, introduced in De
Leeuw (\citeproc{ref-deleeuw_C_77}{1977}). Vector \(e_i\) has \(n\)
elements, with element \(i\) equal to \(+1\), and all other elements
zero. \[
A_{ij}:=(e_i-e_j)(e_i-e_j)'
\] \((e_i-e_j)(e_i-e_j)'\), which means elements \((i,i)\) and \((j,j)\)
are equal to \(+1\), while \((i,j)\) and \((j,i)\) are \(-1\).

\[
\{H_{ij}\}_{st}:=\text{tr}\ Y_s'A_{ij}Y_t
\] Thus \begin{equation}
d_{ij}^2(\alpha)=\text{tr}\ \alpha'A_{ij}\alpha,
\label{eq-anot}
\end{equation}

We also define \begin{equation}
V:=\mathop{\sum}_{1\leq i<j\leq n}w_{ij}A_{ij},
\label{eq-vdef}
\end{equation} and the matrix-valued function \(B\) with
\begin{equation}
B(X):=\mathop{\sum}_{1\leq i<j\leq n}w_{ij}r_{ij}(X)A_{ij},
\label{eq-bdef}
\end{equation} where \begin{equation}
r_{ij}(X):=\begin{cases}\frac{\delta_{ij}}{d_{ij}(X)}&\text{ if }d_{ij}(X)>0,\\
0&\text{ if }d_{ij}(X)=0.\end{cases}\label{eq-rdef}
\end{equation}

If we assume, without loss of generality, that \begin{equation}
\frac12\mathop{\sum}_{1\leq i<j\leq n}w_{ij}\delta_{ij}^2=1,
\label{eq-dnorm}
\end{equation} then \begin{equation}
\sigma(X)=1-\text{tr}\ X'B(X)X+\frac12\text{tr}\ X'VX.
\label{eq-S}
\end{equation}

\(V\) is symmetric with non-positive off-diagonal elements. It is
doubly-centered (rows and columns add up to zero) and thus weakly
diagonally dominant. It follows that it is positive semi-definite (see
Varga (\citeproc{ref-varga_62}{1962}), section 1.5). Because of
irreducibility it has rank \(n-1\), and the vectors in its null space
are all proportional to \(e\), the vector with all elements equal to
\(+1\). The matrix \(B(X)\) is also symmetric, positive semi-definite,
and doubly-centered for each \(X\). It may not be irreducible, because,
for example, \(B(0)=0\).

For the notation in coefficient space we need the square, symmetric,
positive semi-definite matrices \(A_{ij}\) of order \(m\) with elements
\begin{equation}
\{A_{ij}\}_{st}:=\text{tr}\ Y_s'A_{ij}Y_t.\label{eq-atildef}
\end{equation} The squared distance equals \begin{equation}
d_{ij}^2(\alpha):=\alpha'A_{ij}\alpha.\label{eq-dtildef}
\end{equation} Note that we continue to overload symbols such as
\(A_{ij}\) and \(d\). The same is true for \(V\) from \eqref{eq-vdef},
which we define in coefficient space by using definition
\eqref{eq-atildef}, and for \(\tilde B(\alpha)\) from \eqref{eq-bdef}.
In coefficient space \begin{equation}
\sigma(\alpha):=
1-\alpha'B(\alpha)\alpha+\frac12\alpha'V\alpha.\label{eq-stildef}
\end{equation} and the MDS problem is to minimize \eqref{eq-stildef}
over \(\alpha\in\mathbb{R}^m\).

\sectionbreak

\section{SMACOF}\label{smacof}

Define the \emph{Guttman Transform} of configuration \(X\) as
\begin{equation}
\Gamma(X):=V^+B(X)X
\label{eq-guttmanx}
\end{equation} The \emph{SMACOF algorithm} in configuration space is,
with iteration counter \(k\geq 1\), \begin{equation}
X^{(k)}=\Gamma(X^{(k-1)})=\Gamma^k(X^{0}),
\label{eq-smacofx}
\end{equation} where \(\Gamma^0=I\) and
\(\Gamma^k(X):=\Gamma(\Gamma^{k-1}(X))\).

In coefficient space the Guttman Transform is \begin{equation}
\Gamma(\alpha):=V^+B(\alpha)\alpha
\label{eq-guttmana}
\end{equation} The \emph{SMACOF algorithm} in configuration space is
\begin{equation}
\alpha^{(k)}=\Gamma(\alpha^{(k-1)})=\Gamma^k(\alpha^{0}).
\label{eq-smacofx}
\end{equation}

De Leeuw (\citeproc{ref-deleeuw_C_77}{1977}) shows that the SMACOF
iterations \eqref{eq-smacof} tend (in a specific sense described later
in this paper) to a fixed point of the Guttman transform, i.e.~to an
\(X\) with \(\Gamma(X)=X\). If stress is differentiable at \(X\) then
the fixed points of the Guttman transform are stationary points, where
the derivative of stress vanishes. Also \(V^+B(X)X=X\) shows that the
columns of \(X\) are eigenvectors of \(V^+B(X)\), with eigenvalues equal
to one.

The algorithm (without weights) was proposed by Guttman
(\citeproc{ref-guttman_68}{1968}), by differentiating stress and setting
the derivatives equal to zero. This derivation ignores the problem that
stress is not differentiable if one or more distances are zero, and it
also does not say anything about convergence of the algorithm. In De
Leeuw (\citeproc{ref-deleeuw_C_77}{1977}) and De Leeuw and Heiser
(\citeproc{ref-deleeuw_heiser_C_77}{1977}) a new derivation of the
algorithm was given, using a majorization argument based on the
Cauchy-Schwarz inequality. This make it possible to avoid problems with
differentiability, and it leads to a simple convergence proof.

\subsection{Coefficient Space}\label{coefficient-space}

Define the \emph{Guttman Transform} of coefficients \(\alpha\) as
\begin{equation}
\tilde\Gamma(\alpha):=\tilde V^+\tilde B(\alpha)\alpha
\label{eq-guttmana}
\end{equation} The \emph{SMACOF algorithm} is \begin{equation}
\alpha^{(k+1)}=\tilde\Gamma(\alpha^{(k)})=\tilde\Gamma^k(\alpha^{0})
\label{eq-smacofa}
\end{equation}

\sectionbreak

\section{Global Convergence of
SMACOF}\label{global-convergence-of-smacof}

Following De Leeuw (\citeproc{ref-deleeuw_C_77}{1977}) we also define
\begin{align}
\rho(X)&=\mathop{\sum}_{1\leq i<j\leq n} w_{ij}\delta_{ij}d_{ij}(X)=\text{tr}\ X'B(X)X,\\
\eta^2(X)&=\mathop{\sum}_{1\leq i<j\leq n} w_{ij}d_{ij}^2(X)=\text{tr}\ X'VX,
\end{align} and \begin{equation}
\lambda(X)=\frac{\rho(X)}{\eta(X)}.
\end{equation} The main result in De Leeuw
(\citeproc{ref-deleeuw_C_77}{1977}) is:

\begin{enumerate}
\def\labelenumi{\arabic{enumi}.}
\tightlist
\item
  The SMACOF iterates \(X^{(k)}\) are in the compact set
  \(\eta(X)\leq 1\),
\item
  \(\{\rho(X^{k})\}, \{\eta^(X^{k})\}\) and \(\{\lambda(X^{k})\}\) are
  increasing sequences, with limits \(\rho_\infty,\eta_\infty\) and
  \(\lambda_\infty\), where
  \(\eta_\infty=\lambda_\infty=\sqrt{\rho_\infty}\).
\item
  \(\{\sigma(X^{k})\}\) is a decreasing sequence converging to
  \(\sigma_\infty=1-\eta^2_\infty\).
\item
  At any accumulation point \(X_\infty\) of \(\{\sigma(X^{k})\}\) we
  have \(\sigma(X_\infty)=\sigma_\infty\) and
  \(X_\infty=\Gamma(X_\infty)\).
\item
  \(\eta(X^{(k+1)}-X^{(k)})\rightarrow 0\) and thus either
  \(\{\sigma(X^{k})\}\) converges or the set of accumulation points of
  \(\{\sigma(X^{k})\}\) is a continuum.
\end{enumerate}

So, in words, the sequence of stress values converges monotonically. The
sequence of configurations is asymptotically regular and has one or more
accumulation points. Each accumulation point is a fixed point of the
Guttman transform. All accumulation points have the same function value,
which is also the limit of the stress sequence.

But it is important to emphasize that De Leeuw
(\citeproc{ref-deleeuw_C_77}{1977}) did not show that the configuration
sequence \(\{X^{(k)}\}\) actually a Cauchy sequence, and converges to a
configuration \(X_\infty\). This is basically because of the rotational
invariance of the MDS problem. If \(K\) is a rotation matrix,
i.e.~\(K'K=KK'=I\), then \(d_{ij}(XK)=d_{ij}(X)\) for all \(i\) and
\(j\), and thus \(\sigma(XK)=\sigma(X)\). If \(X\) is a fixed point of
the Guttman tranform, then so is \(XK\). Consequently there are no
isolated fixed points, each fixed point is part of a nonlinear continuum
of fixed points.

It is also of interest that De Leeuw
(\citeproc{ref-deleeuw_A_84f}{1984}) showed that if \(X\) is a local
minimizer of stress then \(d_{ij}(X)>0\) for all \(i\) and \(j\) such
that \(w_{ij}\delta_{ij}>0\). Thus, if weights and dissimilarities are
positive, stress is differentiable at a local minimum. De Leeuw
(\citeproc{ref-deleeuw_R_93c}{1993}) shows that stress has only a single
local maximum at \(X=0\), and consequently the only possible stationary
points are local minima and saddle points. It is shown by De Leeuw
(\citeproc{ref-deleeuw_E_19g}{2019}) that all fixed points of the
Guttman transform which are not of full column rank are saddle points.
Thus if \(X\) is a local minimizer of MDS, then adding \(q\) zero
columns to \(X\) produces a saddle point \([X\mid 0]\) for (p+q)MDS. At
saddle points there are always directions in which stress can be
decreased, and consequently the SMACOF algorithm with enough iterations
will eventually get close to a continuum of local minima.

\sectionbreak

\section{Local Convergence of SMACOF}\label{local-convergence-of-smacof}

De Leeuw (\citeproc{ref-deleeuw_A_88b}{1988}) studies the local
convergence of SMACOF, i.e.~the \emph{rate of convergence} of the SMACOF
iterations. There is, however, a problem not adequately addressed in
that article, which has quite a bit of hand-waiving, and some incorrect
statements. If there is no convergence to a single configuration then
the usual rate of convergence is not defined either.

We know the sequence \(\{\eta(X^{(k+1)}-X^{(k)})\}\) tends to zero. We
can measure its rate of convergence by the \emph{ratio factor}
\begin{equation}
q_k=\frac{\eta(X^{(k+1)}-X^{(k)})}{\eta(X^{(k)}-X^{(k-1)})},
\end{equation} and the \emph{root factor} \begin{equation}
r_k=\eta(X^{(k+1)}-X^{(k)})^{1/k}.
\end{equation} There is no guarantee that either \(\{q_k\}\) or
\(\{r_k\}\) converges, but \begin{equation}
\liminf_{k\rightarrow\infty} q_k\leq\liminf_{k\rightarrow\infty} r_k\leq\limsup_{k\rightarrow\infty} r_k\leq\limsup_{k\rightarrow\infty} q_k
\end{equation} Thus if \(\lim_{k\rightarrow\infty} q_k\) exists then
\(\lim_{k\rightarrow\infty} r_k\) exist and the two limit values are
equal. See Rudin (\citeproc{ref-rudin_76}{1976}), p.~68, theorem 3.37.

The rate of convergence of the SMACOF iterations at a fixed point \(X\)
is computed in De Leeuw (\citeproc{ref-deleeuw_A_88b}{1988}) as the
spectral norm \(\kappa(X)=\|\mathcal{D}\Gamma(X)\|_\infty\), i.e.~as the
modulus of the largest eigenvalue of the derivative of the Guttman
transform. There are good reasons to compute this quantity. From
Ostrovski's Theorem (Ostrowski (\citeproc{ref-ostrowski_73}{1973}),
theorem 22.1, p. 151 , Ortega and Rheinboldt
(\citeproc{ref-ortega_rheinboldt_70}{1970}), theorem 10.1.3, p.~300) we
know that if \(\kappa(X)<1\), then \(X\) is a \emph{point of attraction}
of the SMACOF iteration. If we start close enough then the iterations
will converge to \(X\). We also know (Ostrowski
(\citeproc{ref-ostrowski_73}{1973}), theorem 22.2, p.~152) that if
\(\kappa(X)>1\) then \(X\) is a \emph{point of repulsion}, which means
the iterations will diverge when started in a certain solid angle with
\(X\) at its apex.

If we define, with Ortega and Rheinboldt
(\citeproc{ref-ortega_rheinboldt_70}{1970}), chapter 9, for a sequence
\(\{X^{(k)}\}\) converging to \(X\), the ratio and root convergence
factors \begin{equation}
Q_1(\{X^{(k)}\})=\limsup_{k\rightarrow\infty}\frac{\eta(X^{(k+1)}-X)}{\eta(X^{(k)}-X)},
\end{equation} and \begin{equation}
R_1(\{X^{(k)}\})=\limsup_{k\rightarrow\infty}\ \eta(X^{(k)}-X)^{1/k},
\end{equation} then \(\kappa(X)<1\) implies \(R_1(X^{(k)})=\kappa(X)\)
(Ortega and Rheinboldt (\citeproc{ref-ortega_rheinboldt_70}{1970}),
theorem 10.1.4, p.~301). Alo \(R_1(X^{(k)})\) is independent of the norm
we have chosen, which is \(\eta\) in the SMACOF case (Ortega and
Rheinboldt (\citeproc{ref-ortega_rheinboldt_70}{1970}), theorem 9.9.2,
p.~288). Because \(Q_1(X^{(k)})\) depends on the norm we can only assert
that \(Q_1(X^{(k)})\geq R_1(X^{(k)})=\kappa(X)\), although it is
guaranteed in the SMACOF case that some norm exists for which there is
equality (Ortega and Rheinboldt
(\citeproc{ref-ortega_rheinboldt_70}{1970}), Notes and Remarks NR
10.1-5, p.~306).

There is a problem, however, with applying the Ostrowski and
Ortega-Rheinboldt results in our case. Because of the invariance of
distances under rotation we have \(\kappa(X)=1\), no matter what \(X\)
is. And, basically for the same reason, we have not shown that the
SMACOF iterations converge to a fixed point at all. De Leeuw
(\citeproc{ref-deleeuw_A_88b}{1988}) on page 171 acknowledges the
problem. In fact, he proposes a way around the problem in the proof of
theorem 3 on page 175, but the proof is not very convincing, although
convincing enough to get past the reviewers. I will try to be more
specific and precise.

Observe that \begin{equation}
\mathcal{D}\Gamma(X)=V^+\mathcal{D}^2\rho(X),
(\#eq:dgamma)
\end{equation} and \begin{equation}
\mathcal{D}^2\sigma(X)=V-\mathcal{D}^2\rho(X)=V(I-\mathcal{D}\Gamma(X)).
(\#eq:d2sigma)
\end{equation} Because \(\rho\) is convex \(\mathcal{D}^2\rho(X)\) is
symmetric and positive semi-definite. The eigenvalues of
\(\mathcal{D}\Gamma(X)\) are the generalized eigenvalues of the positive
semi-definite matrix pair \((\mathcal{D}^2\rho(X),V)\).

An explicit formula for the derivatives of the Guttman transform, or
equivalently for the second derivatives of stress, was first given by De
Leeuw (\citeproc{ref-deleeuw_A_88b}{1988}). We write
\(\mathcal{D}\Gamma(X)[Y]\) for the derivative evaluated at \(Y\).
Partial derivatives can be obtained by choosing \(Y=e_ie_s'\). Of course
we assume here that \(d_{ij}(X)>0\) for all \(i<j\) with
\(w_{ij}\delta_{ij}>0\). \begin{equation}
\mathcal{D}\Gamma(X)[Y]=V^+\{B(X)Y-H(X,Y)X\},
(\#eq:dgammamat)
\end{equation} with \begin{equation}
H(X,Y)=\mathop{\sum}_{1\leq i<j\leq n}w_{ij}\frac{\delta_{ij}}{d_{ij}(X)}\frac{\text{tr}\ X'A_{ij}Y}{d_{ij}^2(X)}A_{ij}.
(\#eq:hdef)
\end{equation} These formulas, together with equations @ref(eq:dgamma)
and @ref(eq:d2sigma), prove the main results in De Leeuw
(\citeproc{ref-deleeuw_A_88b}{1988}).

\begin{enumerate}
\def\labelenumi{\arabic{enumi}.}
\tightlist
\item
  All eigenvalues of \(\mathcal{D}\Gamma(X)\) are non-negative.
\item
  \textbf{(Homogeneity)} \(X\) is an eigenvector of
  \(\mathcal{D}\Gamma(X)\) with eigenvalue zero.
\item
  \textbf{(Translation)} \(\mathcal{D}\Gamma(X)\) has at least \(p\)
  additional eigenvalues equal to zero, corresponding with eigenvectors
  of the form \(ee_s'\), where \(e\) has all elements equal to one and
  \(e_s\) has one of its elements equal to one and the others equal to
  zero.
\item
  \textbf{(Rotation)} If \(X\) is a fixed point of the Guttman transform
  then \(\mathcal{D}\Gamma(X)\) has at least \(\frac12 p(p-1)\)
  eigenvalues equal to one, corresponding with eigenvectors of the form
  \(XA\), with \(A\) anti-symmetric.
\item
  At a local minimum point \(X\) of stress all eigenvalues of
  \(\mathcal{D}\Gamma(X)\) are less than or equal to one. At a strict
  local minimum they are strictly less than one.
\item
  At a saddle point \(X\) of stress the largest eigenvalue of
  \(\mathcal{D}\Gamma(X)\) is larger than one.
\end{enumerate}

\section{Modified SMACOF}\label{modified-smacof}

To avoid the problems caused by a continuum of fixed points De Leeuw
(\citeproc{ref-deleeuw_A_88b}{1988}) proposes to rotate each SMACOF
iterate to a unique position, for example to principal components. For
any \(X\) with singular value decomposition \(X=K\Lambda L'\) we define
\(\Pi(X)=K\Lambda=XL.\) Note that if one or more singular values are
equal then \(\Pi\) is not unique and becomes a point-to-set map. We can
rotate within each of the spaces corresponding to multiple singular
values, and thus we have not completely eliminated indeterminacy due to
rotation. In the sequel we simply assume that the \(p\) singular values
of \(X\) are different.

Now consider the iterations \begin{equation}
\tilde X^{(k+1)}=\Pi(\Gamma(\tilde X^{(k)}))={(\Pi\Gamma)}^k(\tilde X^{0}).
\end{equation} Let's call this modified SMACOF, or mSMACOF for short.

It looks initialy as if mSMACOF require a great deal of extra
computation, but this is actually not the case. If \(M\) is a rotation
matrix we have \(\Pi(XM)=\Pi(X)\) and \(\Gamma(XM)=\Gamma(X)M\). This
implies that \(\Pi\Gamma\Pi(X)=\Pi\Gamma(X)\) for all \(X\). As a result
the sequences \(\{X^{(k)}\}\) and \(\{\tilde X^{(k)}\}\) have a simple
relationship. If we start the \(\{\tilde X^{(k)}\}\) sequence with
\(X^{0}\) then for each \(k\) we have
\(\sigma(\tilde X^{(k)})=\sigma(X^{(k)})\) and
\(\tilde X^{(k)}=\Pi(X^{(k)})\). Thus the two sequences of stress values
are exactly the same for SMACOF and mSMACOF and the two configurations,
and consequnetly also the accumulation points of the two sequences of
configurations, just differ by a rotation. To get \(\tilde X^{(k)}\) we
can just compute the SMACOF iterate \(X^{(k)}\) and rotate it to
principal components.

The derivative of \(\Pi\Gamma\) is, by the chain rule, \begin{equation}
\mathcal{D}\Pi\Gamma(X)=\mathcal{D}\Pi(\Gamma(X))\mathcal{D}\Gamma(X).
\end{equation} After some computation we find \begin{equation}
\mathcal{D}\Pi(X)[Y]=YL+K\Lambda M,
\end{equation} where \(X=K\Lambda L'\) and \(M\) is the anti-symmetric
matrix with off-diagonal elements \begin{equation}
m_{ij}=-\frac{\lambda_iu_{ij}+\lambda_ju_{ji}}{\lambda_i^2-\lambda_j^2}
\end{equation} with \(U=K'YL\).

If \(X\) satisfies \(\Pi(X)=X\) then \(L=I\) and thus
\(\mathcal{D}\Pi(X)[Y]=Y+XM\) and \(U=K'Y\). We proceed to compute the
eigenvectors and eigenvalues of the Jacobian at an \(X\) with
\(\Pi(X)=X\).

\begin{enumerate}
\def\labelenumi{\arabic{enumi}.}
\tightlist
\item
  If \(Y=XA\) with \(A\) antisymmetric then \(M=-L'AL\) and
  \(\mathcal{D}\Pi(X)[Y]=0\). This corresponds with \(\frac12 p(p-1)\)
  zero eigenvalues
\item
  If \(Y=K\Lambda^{-1}A\) with \(A\) antisymmetric then \(M=0\). This
  gives \(\frac12 p(p-1)\) eigenvectors with eigenvalue one.
\item
  If \(Y=K\Lambda^{-1}D\) wth \(D\) diagonal then \(M=0\). This gives
  another \(p\) eigenvectors with eigenvalue one.
\item
  If \(Y=K_\perp S\) with \(K_\perp\) a basis for the orthogonal
  complement of \(X\) then \(M=0\) and we have another \(p(n-p)\)
  eigenvalues equal to one.
\end{enumerate}

This is a complete set of eigenvectors. It turns out all eigenvalues are
equal to one, except for \(\frac12 p(p-1)\) zero eigenvalues. It follows
that fixed points of \(\Pi\Gamma\) are of the form \(\tilde X=XL\), with
\(X\) a fixed point of \(\Gamma\) and \(L\) the rotation of \(X\) to
principal components). Thus \(\Gamma(\tilde X)=\Gamma(X)L=XL=\tilde X\),
and \(\tilde X\) is a fixed point of both \(\Gamma\) and \(\Pi\Gamma\).
The eigenvalues of \(\mathcal{D}\Pi\Gamma(\tilde X)\) are the same as
the eigenvalues of \(\mathcal{D}\Gamma(X)\), except for the
\(\frac12 p(p-1)\) eigenvalues equal to one, corresponding with \(XA\)
for antisymmetric A, which are replaced by zero eigenvalues. The
Ostrowski and Ortega-Rheinboldt results now apply directly to
\(\Pi\Gamma\). If there are no zero distances, if the singular values of
\(X\) are different, and if the largest eigenvalue of
\(\mathcal{D}\Pi\Gamma(X)\) is less than one, then \(X\) is a point of
attraction, the SMACOF iterations converge to \(X\) if started
sufficiently close to it (in the \emph{attraction ball} of \(X\)), and
\(R_1(X)\) is equal to the largest eigenvalue.

\section{Software}\label{software}

The appendix has an ad-hoc version of the \texttt{smacof} function. Its
arguments are

\begin{Shaded}
\begin{Highlighting}[]
\FunctionTok{args}\NormalTok{(smacof)}
\end{Highlighting}
\end{Shaded}

\begin{verbatim}
function (delta, w, p = 2, xold = torgerson(delta, p), pca = FALSE, 
    verbose = FALSE, eps = 1e-15, itmax = 10000) 
NULL
\end{verbatim}

The function expects both dissimilarities and weights to be of class
\texttt{dist}. The default initial configuration uses classical MDS
(Torgerson (\citeproc{ref-torgerson_58}{1958})). We iterate until either
the distance \(\eta(X^{(k)}-X^{(k-1)})\) between successive
configurations is less than \texttt{eps}, which has the ridiculously
small default value of \ensuremath{10^{-15}}, or until the maximum
number of iterations \texttt{itmax} is reached, which defaults to the
ridiculously large default value of 10000. If \texttt{pca} is TRUE all
iterates are rotated to principal components. If \texttt{verbose} is
TRUE we print, for each iteration, the iteration number \(k\), the
stress \(\sigma(X^{(k)})\), the change \(\eta(X^{(k)}-X^{(k-1)})\), and
the root and ratio convergence factors.

The second file \texttt{derivative.R} computes the Jacobian of the
iteration functions at the limit (i.e.~at the point where SMACOF stops).
There are actually six functions in the file. The three functions
\texttt{dGammaA,\ dPiA} and \texttt{dPiGammaA} use the analytical
expressions we have derived earlier for
\(\mathcal{D}\Gamma(X), \mathcal{D}\Pi(X)\) and
\(\mathcal{D}\Pi\Gamma(X).\) The corresponding functions
\texttt{dGammaN,\ dPiN} and \texttt{dPiGammaN} numerically compute the
Jacobian using the \texttt{jacobian} function from the \texttt{numDeriv}
package (Gilbert and Varadhan
(\citeproc{ref-gilbert_varadhan_19}{2019})). They are basically used to
check the formulas.

\section{Examples}\label{examples}

\subsection{Ekman}\label{ekman}

Time for a numerical example. We will use the color similarity data of
Ekman (\citeproc{ref-ekman_54}{1954}), taken from the SMACOF package (De
Leeuw and Mair (\citeproc{ref-deleeuw_mair_A_09c}{2009})). We transform
Ekman's similarities \(s_{ij}\) to dissimilarities \(\delta_{ij}\) using
\(\delta_{ij}=(1-s_{ij})^3\), and we use unit weights in \(W\).

SMACOF requires 52 iterations for convergence. The minimum of stress is
0.0110248119, the root factor is equal to 0.5085800483 and the ratio
factor is 0.5329246268. The eigenvalues of \(\mathcal{D}\Gamma(X)\) are

\begin{verbatim}
0.999999999999999 
0.538510668196408 
0.532498554224166 
0.529669334191043 
0.525541097754551 
0.519602718218807 
0.516533687841054 
0.513727908691608 
0.510295310409166 
0.506357695934493 
0.495649713446001 
0.481518780073304 
0.469548625536451 
0.445038589020916 
0.410479252654484 
0.394284315478173 
0.357670750114585 
0.348395413045201 
0.330341477528788 
0.313520384427863 
0.287598015841956 
0.280399104121623 
0.267278474596759 
0.247631495462991 
0.216576009083047 
0.000000000000000 
0.000000000000000 
0.000000000000000 
\end{verbatim}

As an aside, the Ekman example (with this particular transformation of
the similarities) is special, because the two-dimensional solution is
actually the global minimum over all configurations (i.e.~the global
mnimum of MDS for all \(1\leq p\leq n\)). As shown in De Leeuw
(\citeproc{ref-deleeuw_U_14b}{2014}) this follows from the fact that the
two unit eigenvalues of \(V^+B(X)\) are actually its two largest
eigenvalues, and consequently \(X\) is the solution of the convex
full-dimensional nMDS relaxation of the MDS problem (De Leeuw et al.
(\citeproc{ref-deleeuw_groenen_mair_E_16e}{2016})). The eigenvalues of
\(V^+B(X)\) are

\begin{verbatim}
1.000000000000000 
1.000000000000000 
0.923497086367287 
0.907901212970889 
0.862936584878895 
0.852692003103344 
0.829803620857357 
0.814556167662381 
0.793238576338502 
0.791651722456649 
0.786442678118770 
0.747679475705025 
0.728268247434340 
0.000000000000001 
\end{verbatim}

We now set \texttt{pca=TRUE} and run mSMACOF. It requires 51 iterations
for convergence. The minimum of stress is 0.0110248119, and the root
factor is equal to 0.5080118475 and the ratio factor is 0.5718672565.
The eigenvalues of \(\mathcal{D}\Pi\Gamma(X)\) at the fixed point \(X\)
are

\begin{verbatim}
0.538510668196407 
0.532498554224166 
0.529669334191045 
0.525541097754552 
0.519602718218807 
0.516533687841055 
0.513727908691609 
0.510295310409168 
0.506357695934493 
0.495649713446002 
0.481518780073306 
0.469548625536451 
0.445038589020917 
0.410479252654486 
0.394284315478174 
0.357670750114585 
0.348395413045201 
0.330341477528788 
0.313520384427864 
0.287598015841957 
0.280399104121623 
0.267278474596758 
0.247631495462991 
0.216576009083047 
0.000000000000000 
0.000000000000000 
0.000000000000000 
0.000000000000000 
\end{verbatim}

They are equal to the eigenvalues of \(\mathcal{D}\Gamma(X)\) and the
root convergence factor \(R_1(\{X^{(k)}\})\) is 0.5385106682.

\subsection{De Gruijter}\label{de-gruijter}

The second example are dissimilarties between nine Dutch political
parties, collected a long time ago by De Gruijter
(\citeproc{ref-degruijter_67}{1967}).

\begin{verbatim}
      KVP PvdA  VVD  ARP  CHU  CPN  PSP   BP
PvdA 2.63                                   
VVD  2.27 3.72                              
ARP  1.60 2.64 2.46                         
CHU  1.80 3.22 1.97 0.20                    
CPN  4.54 2.12 5.13 4.84 4.80               
PSP  3.73 1.59 4.55 3.73 4.08 1.08          
BP   4.18 4.22 3.90 4.28 3.96 3.34 3.88     
D66  3.17 2.47 1.67 3.13 3.04 4.42 3.36 4.36
\end{verbatim}

We analyze the data with SMACOF using three dimensions.

SMACOF requires 779 iterations for convergence. The minimum of stress is
0.003442194, the root factor is equal to 0.9565946372, and the ratio
factor is 0.9684304863. The eigenvalues of \(\mathcal{D}\Gamma(X)\) are

\begin{verbatim}
1.000000000000001 
1.000000000000000 
1.000000000000000 
0.965505429805660 
0.940592046981168 
0.919047686446446 
0.863993920924431 
0.822263696226349 
0.810020971414307 
0.771947328045071 
0.728539213663602 
0.712621684571535 
0.659318624263222 
0.638485365688918 
0.624552464148482 
0.616085432872672 
0.557480497708779 
0.501954186928526 
0.458630667850014 
0.450598367846591 
0.355243400669833 
0.309186479174062 
0.247708397109091 
0.000000000000002 
0.000000000000002 
0.000000000000000 
0.000000000000000 
\end{verbatim}

In this case there is no guarantee that we have the three-dimensional
global minimum. The unit eigenvalues of the matrix \(V^+B(X)\) are not
the three largest ones.

\begin{verbatim}
1.079524009371955 
1.032606649163671 
1.000000000000000 
1.000000000000000 
1.000000000000000 
0.986706272372898 
0.971839080877692 
0.906211919383164 
0.000000000000001 
\end{verbatim}

mSMACOF in three dimensions requires 778 iterations for convergence. The
minimum of stress is 0.003442194, the root factor is equal to
0.9565709099, and the ratio factor is 0.9714170285. The eigenvalues of
\(\mathcal{D}\Pi\Gamma(X)\) are

\begin{verbatim}
0.965505429805661 
0.940592046981168 
0.919047686446446 
0.863993920924433 
0.822263696226351 
0.810020971414307 
0.771947328045072 
0.728539213663603 
0.712621684571534 
0.659318624263222 
0.638485365688918 
0.624552464148483 
0.616085432872671 
0.557480497708778 
0.501954186928526 
0.458630667850015 
0.450598367846592 
0.355243400669831 
0.309186479174061 
0.247708397109091 
0.000000000000001 
0.000000000000001 
0.000000000000001 
0.000000000000000 
0.000000000000000 
0.000000000000000 
0.000000000000000 
\end{verbatim}

Again, as expected, the eigenvalues for SMACOF and mSMACOF are the same
and the root convergence factor is the spectral norm of the mSMACOF
derivative, which is 0.9655054298.

\section{Conclusion}\label{conclusion}

From a practical point of view there is no need to ever use modified
SMACOF. We only use the Guttman transform, and compute the largest
eigenvalue of its derivative at the solution \(X\), ignoring the
\(\frac12 p(p-1)\) trivial unit eigenvalues. That largest eigenvalue, if
it is strictly smaller than one, is the root convergence factor. In that
case SMACOF can be said to converge to \(\Pi(X)\), which is the SMACOF
solution rotated to principal components.

\sectionbreak

\section*{References}\label{references}
\addcontentsline{toc}{section}{References}

\protect\phantomsection\label{refs}
\begin{CSLReferences}{1}{1}
\bibitem[\citeproctext]{ref-bentler_weeks_78}
Bentler, P. M., and D. G. Weeks. 1978. {``Restricted Multidimensional
Scaling Models.''} \emph{Journal of Mathematical Psychology} 17:
138--51.

\bibitem[\citeproctext]{ref-degruijter_67}
De Gruijter, Dato N. M. 1967. \emph{{The Cognitive Structure of Dutch
Political Parties in 1966}}. Report Nos. E019-67. Psychological
Institute, University of Leiden.

\bibitem[\citeproctext]{ref-deleeuw_R_93c}
De Leeuw, J. 1993. \emph{Fitting Distances by Least Squares}. Preprint
Series No. 130. UCLA Department of Statistics.
\url{https://jansweb.netlify.app/publication/deleeuw-r-93-c/deleeuw-r-93-c.pdf}.

\bibitem[\citeproctext]{ref-deleeuw_U_14b}
De Leeuw, J. 2014. {``{Bounding, and Sometimes Finding, the Global
Minimum in Multidimensional Scaling}.''} UCLA Department of Statistics.
\url{https://jansweb.netlify.app/publication/deleeuw-u-14-b/deleeuw-u-14-b.pdf}.

\bibitem[\citeproctext]{ref-deleeuw_E_19g}
De Leeuw, J. 2019. {``{Fitting Distances by Least Squares}.''}
\url{https://jansweb.netlify.app/publication/deleeuw-r-93-c/deleeuw-r-93-c.pdf}.

\bibitem[\citeproctext]{ref-deleeuw_C_77}
De Leeuw, Jan. 1977. {``Applications of Convex Analysis to
Multidimensional Scaling.''} In \emph{Recent Developments in
Statistics}, edited by J. R. Barra, F. Brodeau, G. Romier, and B. Van
Cutsem. North Holland Publishing Company.

\bibitem[\citeproctext]{ref-deleeuw_A_84f}
De Leeuw, Jan. 1984. {``{Differentiability of Kruskal's Stress at a
Local Minimum}.''} \emph{Psychometrika} 49: 111--13.

\bibitem[\citeproctext]{ref-deleeuw_A_88b}
De Leeuw, Jan. 1988. {``Convergence of the Majorization Method for
Multidimensional Scaling.''} \emph{Journal of Classification} 5:
163--80.

\bibitem[\citeproctext]{ref-deleeuw_groenen_mair_E_16e}
De Leeuw, Jan, Patrick J. F. Groenen, and Patrick Mair. 2016.
{``Full-Dimensional Scaling.''}
\url{https://jansweb.netlify.app/publication/deleeuw-groenen-mair-e-16-e/deleeuw-groenen-mair-e-16-e.pdf}.

\bibitem[\citeproctext]{ref-deleeuw_heiser_C_80}
De Leeuw, Jan, and Willem J. Heiser. 1980. {``Multidimensional Scaling
with Restrictions on the Configuration.''} In \emph{Multivariate
Analysis, Volume {V}}, edited by P. R. Krishnaiah. North Holland
Publishing Company.

\bibitem[\citeproctext]{ref-deleeuw_mair_A_09c}
De Leeuw, Jan, and Patrick Mair. 2009. {``{Multidimensional Scaling
Using Majorization: SMACOF in R}.''} \emph{Journal of Statistical
Software} 31 (3): 1--30.
\url{https://www.jstatsoft.org/article/view/v031i03}.

\bibitem[\citeproctext]{ref-deleeuw_heiser_C_77}
De Leeuw, J., and W. J. Heiser. 1977. {``Convergence of Correction
Matrix Algorithms for Multidimensional Scaling.''} Chap. 32 in
\emph{Geometric Representations of Relational Data}, edited by J. C.
Lingoes. Mathesis Press.

\bibitem[\citeproctext]{ref-ekman_54}
Ekman, Gosta. 1954. {``{Dimensions of Color Vision}.''} \emph{Journal of
Psychology} 38: 467--74.

\bibitem[\citeproctext]{ref-gilbert_varadhan_19}
Gilbert, P., and R. Varadhan. 2019. \emph{{numDeriv: Accurate Numerical
Derivatives}}. \url{https://CRAN.R-project.org/package=numDeriv}.

\bibitem[\citeproctext]{ref-guttman_68}
Guttman, Louis. 1968. {``{A General Nonmetric Technique for Fitting the
Smallest Coordinate Space for a Configuration of Points}.''}
\emph{Psychometrika} 33: 469--506.

\bibitem[\citeproctext]{ref-kruskal_64a}
Kruskal, Joseph B. 1964a. {``{Multidimensional Scaling by Optimizing
Goodness of Fit to a Nonmetric Hypothesis}.''} \emph{Psychometrika} 29:
1--27.

\bibitem[\citeproctext]{ref-kruskal_64b}
Kruskal, Joseph B. 1964b. {``{Nonmetric Multidimensional Scaling: a
Numerical Method}.''} \emph{Psychometrika} 29: 115--29.

\bibitem[\citeproctext]{ref-ortega_rheinboldt_70}
Ortega, J. M., and W. C. Rheinboldt. 1970. \emph{{Iterative Solution of
Nonlinear Equations in Several Variables}}. Academic Press.

\bibitem[\citeproctext]{ref-ostrowski_73}
Ostrowski, A. M. 1973. \emph{Solution of Equations in Euclidean and
Banach Spaces}. Third Edition of Solution of Equations and Systems of
Equations. Academic Press.

\bibitem[\citeproctext]{ref-rudin_76}
Rudin, W. 1976. \emph{Principles of Mathematical Analysis}. Third
Edition. McGraw-Hill.

\bibitem[\citeproctext]{ref-torgerson_58}
Torgerson, Warren S. 1958. \emph{{Theory and Methods of Scaling}}.
Wiley.

\bibitem[\citeproctext]{ref-varga_62}
Varga, R. S. 1962. \emph{{Matrix Iterative Analysis}}. Prentice Hall.

\end{CSLReferences}




\end{document}
